\section{Setting, hypotheses and exact identities}\label{sec:setting}

\subsection{Runs, preparation and the occupation count}\label{sec:runs}

The walk takes $N$ steps and each step has a sign, $\plus$ or $\minus$. We call
the positions $1,\dots,N$ slots. The signs come in maximal runs of equal sign,
and consecutive runs have opposite signs. All runs after the first are
independent. A run of sign $\sgn\in\{\plus,\minus\}$ has a length distributed as
$X_\sgn$, a random variable on $\{1,2,\dots\}$ with mass function
$f_\sgn(n)=\Prob(X_\sgn=n)$ and survival function
\[
 T_\sgn(x)=\Prob(X_\sgn\ge x),\qquad x\in\R .
\]
We extend $f_\sgn$ by zero to nonpositive integers, and we write $\osgn$ for the
sign opposite to $\sgn$.

The first run is the preparation. It has a sign $\Sigma$ and a length $D$, and
the pair $(\Sigma,D)$ is independent of all later runs. We put
\[
 \nu_\sgn(d)=\Prob(\Sigma=\sgn,\,D=d),\qquad
 Q_\sgn(N)=\Prob(\Sigma=\sgn,\,D\ge N),
\]
and we assume that $D$ is finite almost surely. A fresh start of sign $\sgn_0$
is the preparation $\nu_{\sgn_0}=f_{\sgn_0}$ and $\nu_{\osgn_0}=0$. The
stationary preparation has $\nu_\sgn=T_\sgn/s$ with $s$ as in~\eqref{eq:means}
below. It is the law of the current run seen from a typical time of a
two-sided stationary sequence of runs, and it is the equilibrated start of the
physics literature.

Let $A_N$ be the number of plus slots among $1,\dots,N$, and
\[
 P_N(k)=\Prob(A_N=k),\qquad 0\le k\le N .
\]
The walk that moves right on a plus slot and left on a minus slot is at
$2A_N-N$ after $N$ steps. For a fresh start of sign $\sgn$ we
write $F^\sgn_N(k)$ for $P_N(k)$, with $F^\sgn_N(k)=0$ for $k\notin\{0,\dots,N\}$.
We also need the law of the count of the starting sign,
\[
 \Pi^\plus_m(b)=F^\plus_m(b),\qquad \Pi^\minus_m(b)=F^\minus_m(m-b),
\]
so $\Pi^\sgn_m(b)$ is the probability that a fresh walk of sign $\sgn$ has
exactly $b$ slots of its own sign among its first $m$ slots.

\subsection{Hypotheses}\label{sec:hypotheses}

The hypotheses carry fixed names, and every statement says which it uses.
\begin{itemize}[leftmargin=4.5em,labelsep=0.8em]
\item[(RV)] For each sign, $T_\sgn$ is regularly varying at infinity with index
$-\alpha_\sgn$, where $1<\alpha_\sgn<2$. We write $\alpha_*=\min(\alpha_\plus,\alpha_\minus)$.
\item[(LR)] For each sign there is $\alpha_\sgn\in(1,2)$ with
$f_\sgn(n)\sim\alpha_\sgn T_\sgn(n)/n$ as $n\to\infty$.
\item[(SPAN)] For each sign, the support of $X_\sgn$ is not contained in any
arithmetic progression $b+h\Z$ with $h\ge2$.
\end{itemize}
A hypothesis assumed ``for the plus sign'' is assumed for that sign only. Under
(LR) we have $T_\sgn(n+1)/T_\sgn(n)=1-(\alpha_\sgn+o(1))/n$, so the representation
theorem~\cite[Theorem~1.3.1]{bgt1989} shows that (LR) implies (RV) with the same
index. (LR) says that the mass function is regularly varying, not only the
tail. (SPAN) is the aperiodicity needed for local limit theorems. Under (RV) the
runs have finite means, and we put
\begin{equation}\label{eq:means}
 \mu_\sgn=\E X_\sgn,\qquad s=\mu_\plus+\mu_\minus,\qquad p=\mu_\plus/s ,
\end{equation}
so $p$ is the long-run fraction of plus slots and $A_N/N\to p$ in probability.

\subsection{Scales and the moving profile}\label{sec:scales}

Under (RV) choose for each sign an increasing function
$a_\sgn:(0,\infty)\to(0,\infty)$ with $t\,T_\sgn(a_\sgn(t))\to1$ as
$t\to\infty$. So $a_\sgn$ is regularly varying with index $1/\alpha_\sgn$, and
$a_\sgn(m)$ is the scale of the stable fluctuation of a sum of $m$ runs of sign
$\sgn$. Put
\begin{equation}\label{eq:scales}
 b_{\sgn,N}=a_\sgn(N/s),\qquad B_N=\max(b_{\plus,N},b_{\minus,N}),\qquad
 \theta_{\sgn,N}=b_{\sgn,N}/B_N\in(0,1],
\end{equation}
and $x_{N,k}=(k-pN)/B_N$. The walk makes about $N/s$ runs of each sign in $N$
slots, so $b_{\sgn,N}$ is the fluctuation scale of the total length of the runs
of sign $\sgn$, and $B_N$ is the larger of the two. When
$\alpha_\plus\ne\alpha_\minus$ one ratio $\theta_{\sgn,N}$ tends to zero, and when
the indices are equal the ratios may oscillate, so we keep them in the
statements. We collect the facts about regular variation that recur.

\begin{lemma}[regular variation]\label{lem:rv}
Let $f$ be regularly varying with index $\varrho$.
\begin{enumerate}[(i)]
\item (Uniform convergence.) The supremum of $|f(x)/f(y)-1|$ over $x,y\ge x_0$
with $|x/y-1|\le r$ tends to zero as $x_0\to\infty$ and
$r\to0$~\cite[Theorems~1.2.1 and~1.5.2]{bgt1989}.
\item (Potter's bounds~\cite[Theorem~1.5.6]{bgt1989}.) For every $\delta>0$ there is
$x_\delta$ with $f(y)/f(x)\le2\max\{(y/x)^{\varrho+\delta},(y/x)^{\varrho-\delta}\}$ for
$x,y\ge x_\delta$. Under (RV), for every $c\ge1$ there is $C_c$ with
$T_\sgn(x/c)\le C_cT_\sgn(x)$ for $x\ge c$.
\item \cite[Theorem~1.5.3]{bgt1989} If $\varrho>0$ and $f$ is bounded above and away
from zero on compact subsets of $[1,\infty)$, then $\sup_{1\le y\le x}f(y)\le C_ff(x)$
for $x\ge1$.
\item (Karamata~\cite[Propositions~1.5.8 and~1.5.10]{bgt1989}.) Under (RV) there is
$C_{\mathrm K}$ such that, for both signs and all $w\ge1$,
\begin{equation}\label{eq:karamata}
 \sum_{n<w}n\,T_\sgn(n)\le C_{\mathrm K}\,w^2T_\sgn(w),\qquad
 \sum_{n\ge w}T_\sgn(n)\le C_{\mathrm K}\,wT_\sgn(w),
\end{equation}
and $\sum_{n\ge w}T_\sgn(n)\sim wT_\sgn(w)/(\alpha_\sgn-1)$ as $w\to\infty$.
\item Under (RV), $B_N$ is nondecreasing and regularly varying with index
$1/\alpha_*$. Moreover $B_NT_\sgn(N)\to0$, $NT_\sgn(b_{\sgn,N})\to s$ and
$B_N\log N=o(N)$.
\item Under (RV), for both signs, all large $N$ and all $L\ge b_{\sgn,N}$,
\begin{equation}\label{eq:NT-bounds}
 \frac s4\Bigl(\frac L{b_{\sgn,N}}\Bigr)^{-\alpha_\sgn-1}\le NT_\sgn(L)\le4s\Bigl(\frac L{b_{\sgn,N}}\Bigr)^{-(\alpha_\sgn+1)/2}.
\end{equation}
\end{enumerate}
\end{lemma}

\begin{proof}
Parts (i) to (iii) are the cited theorems, and $T_\sgn(x/c)\le C_cT_\sgn(x)$ follows
from (i) for large $x$ and from the positivity of the nonincreasing $T_\sgn$ for the
others. For (iv), $nT_\sgn(n)$ has index $1-\alpha_\sgn>-1$ and $T_\sgn$ has index
$-\alpha_\sgn<-1$, so the cited propositions give the asymptotics, and positivity gives
the bounds for all $w\ge1$. For (v), each $a_\sgn$ is increasing, so $B_N$ is
nondecreasing. If $\alpha_\plus\ne\alpha_\minus$, the ratio $b_{\plus,N}/b_{\minus,N}$
is regularly varying with a nonzero index, so $B_N=b_{\sgn,N}$ for the sign with the
smaller index and all large $N$. If the indices are equal, $B_{cN}/B_N$ lies between
the two ratios $b_{\sgn,cN}/b_{\sgn,N}$, which tend to $c^{1/\alpha_*}$. Then
$B_NT_\sgn(N)$ has index $1/\alpha_*-\alpha_\sgn<0$, and a regularly varying function
with a negative index tends to zero~\cite[Proposition~1.5.1]{bgt1989}. The definition of
$a_\sgn$ with $t=N/s$ gives $NT_\sgn(b_{\sgn,N})\to s$, and $B_N\log N=o(N)$ because
$1/\alpha_*<1$. For (vi), take $\delta=(\alpha_\sgn-1)/2\le\frac12$ in (ii). For large
$N$ we have $b_{\sgn,N}\ge x_\delta$ and $s/2\le NT_\sgn(b_{\sgn,N})\le2s$, and
$(L/b_{\sgn,N})^{-\alpha_\sgn-\delta}\ge(L/b_{\sgn,N})^{-\alpha_\sgn-1}$ for $L\ge b_{\sgn,N}$.
\end{proof}

Let $K_\alpha=-\Gamma(1-\alpha)>0$ for $1<\alpha<2$, and let $Z_\plus$ and
$Z_\minus$ be independent with
\[
 \E e^{itZ_\sgn}=\exp\bigl(K_{\alpha_\sgn}(-it)^{\alpha_\sgn}\bigr),
\]
with the principal branch. These are the totally skewed stable laws that arise
as limits of centered sums of the runs, and we write $g_\sgn$ for the density
of $Z_\sgn$. For $\theta=(\theta_\plus,\theta_\minus)\in[0,1]^2$ let $h_\theta$
be the density of
\begin{equation}\label{eq:profile}
 W_\theta=(1-p)\,\theta_\plus Z_\plus-p\,\theta_\minus Z_\minus ,
\end{equation}
let $\Theta$ be the compact set of $\theta\in[0,1]^2$ with
$\max(\theta_\plus,\theta_\minus)=1$, and put $\theta_N=(\theta_{\plus,N},\theta_{\minus,N})\in\Theta$.
The moving profile is $h_N=h_{\theta_N}$. It is a stable density when
$\alpha_\plus=\alpha_\minus$, and otherwise a convolution of two stable densities
with different indices. We write $H_N=\sup_xh_N(x)$ for its height. For each
$N$ let $\tau=\tau_N$ be a sign with $b_{\tau,N}=B_N$, taking $\tau=\plus$ if
the two scales are equal. We call $\tau$ the dominant sign and $\bar\tau$ the
other sign, and we put $\varsigma_\plus=1$ and $\varsigma_\minus=-1$.

\begin{lemma}[the profile]\label{lem:profile}
Assume (RV).
\begin{enumerate}[(a)]
\item There are $C_1,C_2$ with $\sup_NH_N\le C_1$ and $|h_N(x)-h_N(y)|\le C_2|x-y|$ for
all $N$ and $x,y$.
\item $\varepsilon_h(R)=\sup_N\sup_{|x|\ge R}h_N(x)$ and
$\varepsilon'_h(R)=\sup_N\int_{|x|>R}h_N(x)\,dx$ tend to zero as $R\to\infty$.
\item The map $\theta\mapsto h_\theta$ is continuous from $\Theta$ into the bounded
continuous functions with the uniform norm. So $h_N\to h_\theta$ uniformly along every
subsequence with $\theta_N\to\theta$, and $\lVert h_{N'}-h_N\rVert_\infty\to0$ whenever
$N'/N\to1$.
\item $\inf_NH_N>0$, and $\inf_N\min_{|x|\le K}h_N(x)>0$ for every $K\ge0$.
\item For every $N$ and $x\in\R$,
\begin{equation}\label{eq:bulk-profile}
 h_N(x)=\frac{s}{\mu_{\bar\tau}}\int_\R g_{\bar\tau}(y)\,
 g_\tau\Bigl(\varsigma_\tau\frac{sx}{\mu_{\bar\tau}}+\frac{\mu_\tau}{\mu_{\bar\tau}}\theta_{\bar\tau,N}\,y\Bigr)dy .
\end{equation}
\end{enumerate}
\end{lemma}

\begin{proof}
\emph{Bounds and continuity.} Under (RV) the means are finite, so $0<p<1$. For
$1<\alpha<2$ and the principal branch, $(\mp it)^\alpha=|t|^\alpha e^{\mp i\pi\alpha\,\mathrm{sgn}(t)/2}$, so
the real parts of $(-it)^\alpha$ and $(it)^\alpha$ equal $|t|^\alpha\cos\frac{\pi\alpha}2<0$, and
$(-ict)^\alpha=c^\alpha(-it)^\alpha$ for $c\ge0$. Hence $W_\theta$ has the
characteristic function
\[
 \widehat h_\theta(t)=\exp\Bigl(K_{\alpha_\plus}\bigl((1-p)\theta_\plus\bigr)^{\alpha_\plus}(-it)^{\alpha_\plus}
 +K_{\alpha_\minus}\bigl(p\,\theta_\minus\bigr)^{\alpha_\minus}(it)^{\alpha_\minus}\Bigr),
\]
which is continuous in $\theta$ for each $t$. With
$c_\sgn=K_{\alpha_\sgn}\lvert\cos\frac{\pi\alpha_\sgn}2\rvert$ and
$c=\min(c_\plus(1-p)^{\alpha_\plus},c_\minus p^{\alpha_\minus})>0$, one coordinate of
$\theta\in\Theta$ equals $1$, so $|\widehat h_\theta(t)|\le\Phi(t)$, where $\Phi(t)=1$ for
$|t|<1$ and $\Phi(t)=e^{-c|t|^{\alpha_*}}$ for $|t|\ge1$. Both $\Phi$ and $|t|\Phi(t)$
are integrable. Fourier inversion gives the continuous density
$h_\theta(x)=\frac1{2\pi}\int e^{-itx}\widehat h_\theta(t)\,dt$, with
$\lVert h_\theta\rVert_\infty\le C_1=\frac1{2\pi}\int\Phi$ and Lipschitz constant
$C_2=\frac1{2\pi}\int|t|\Phi(t)\,dt$, and $h_\theta$ vanishes at infinity by the
Riemann--Lebesgue lemma. This gives (a). Also
$\lVert h_\theta-h_{\theta'}\rVert_\infty\le\frac1{2\pi}\int|\widehat h_\theta-\widehat h_{\theta'}|$,
and dominated convergence with the bound $2\Phi$ gives the continuity in (c). If
$N'/N\to1$, then $a_\sgn(N'/s)/a_\sgn(N/s)\to1$ by Lemma~\ref{lem:rv}(i), so
$\theta_{N'}-\theta_N\to0$, and the uniform continuity of $\theta\mapsto h_\theta$ on the
compact $\Theta$ gives the last claim of (c).

\emph{Tails.} The set $\{h_\theta:\theta\in\Theta\}$ is compact. Given $\varepsilon>0$
take a finite $\varepsilon$-net $h_{\theta_1},\dots,h_{\theta_q}$ of it and $R$ with
$h_{\theta_i}<\varepsilon$ on $|x|\ge R$ for all $i$. Then $h_N<2\varepsilon$ on $|x|\ge R$
for every $N$, so $\varepsilon_h(R)\to0$. If $\varepsilon'_h(R)\not\to0$, there are
$\varepsilon>0$, $R_i\to\infty$ and $N_i$ with $\int_{|x|>R_i}h_{N_i}\ge\varepsilon$, and by
compactness we may assume $\theta_{N_i}\to\theta$. Then $h_{N_i}\to h_\theta$ uniformly,
so in $L^1$ by Scheff\'e's lemma, and
$\int_{|x|>R_i}h_{N_i}\le\lVert h_{N_i}-h_\theta\rVert_1+\int_{|x|>R_i}h_\theta\to0$, a
contradiction. This proves (b).

\emph{Positivity.} Fix a sign, write $Z=Z_\sgn$, $g=g_\sgn$ and $\alpha=\alpha_\sgn$, and let
$Z_1,Z_2,\dots$ be independent copies of $Z$. The scaling of the characteristic
function shows that $Z_1+\dots+Z_n$ has the law of $n^{1/\alpha}Z$. If $Z\ge0$ almost
surely, then $\delta=\E\min(Z,1)>0$, and by the strong law of large numbers
$n^{-1/\alpha}(Z_1+\dots+Z_n)\ge n^{1-1/\alpha}\cdot\frac1n\sum_{i\le n}\min(Z_i,1)\to\infty$
almost surely, which contradicts the tightness of a sequence with the law of $Z$.
The same holds for $-Z$. So the open set $U=\{g>0\}$ meets both half lines. The
density $g^{*n}$ of $Z_1+\dots+Z_n$ is continuous and positive at $y_1+\dots+y_n$
whenever $y_1,\dots,y_n\in U$, and it equals $n^{-1/\alpha}g(n^{-1/\alpha}\,\cdot)$. Hence
$y_1,\dots,y_n\in U$ implies $n^{-1/\alpha}(y_1+\dots+y_n)\in U$. Take
$I=(a-\varepsilon,a+\varepsilon)\subset U$ with $0<\varepsilon<a$ and put $c_n=n^{1-1/\alpha}$.
Then $c_nI\subset U$ for every $n$. Since $c_n\to\infty$ and $c_{n+1}/c_n\to1$, the
intervals $c_nI$ with $n\ge n_0$ overlap in turn, so $U$ contains a half line
$(x_0,\infty)$. In the same way $U$ contains $(-\infty,-x_0)$. Every real number is the
sum of a point of each half line, so the case $n=2$ gives $U=\R$. Now let $\theta\in\Theta$
and $\theta_\sgn=1$. Then $W_\theta=aZ_\sgn+Y$ with $a\in\{1-p,-p\}$ and $Y$ independent
of $Z_\sgn$, so $x\mapsto\E[|a|^{-1}g_\sgn((x-Y)/a)]$ is a continuous positive density of
$W_\theta$, and it equals $h_\theta$. The maps $\theta\mapsto\sup_xh_\theta(x)$ and
$\theta\mapsto\min_{|x|\le K}h_\theta(x)$ are continuous and positive on the compact
$\Theta$, so their minima are positive. This proves (d).

\emph{The formula.} Let $\tau=\plus$. Given $Z_\minus=y$, the variable
$(1-p)Z_\plus-p\theta_{\minus,N}Z_\minus$ has the density
$(1-p)^{-1}g_\plus((x+p\theta_{\minus,N}y)/(1-p))$ at $x$. Since $1-p=\mu_\minus/s$ and
$p/(1-p)=\mu_\plus/\mu_\minus$, integrating against $g_\minus(y)\,dy$
gives~\eqref{eq:bulk-profile}. For $\tau=\minus$ we condition on $Z_\plus$.
\end{proof}

\subsection{Renewal kernels}\label{sec:kernels}

For each sign let $u_\sgn$ be the renewal mass function of the runs of sign $\sgn$,
\[
 u_\sgn(0)=1,\qquad u_\sgn(j)=\sum_{m\ge1}f_\sgn^{*m}(j)\quad(j\ge1),
\]
the probability that some partial sum of independent copies of $X_\sgn$ equals
$j$. We also need
\[
 v_\sgn(j)=\sum_{m\ge1}m\,f_\sgn^{*m}(j),
\]
and the iterated kernels $u^\sgn_m=f_\sgn^{*m}$ and $q^\sgn_m=T_\sgn*u^\sgn_{m-1}$
for $m\ge1$, with $u^\sgn_0=\delta_0$. Let $X^\sgn_1,X^\sgn_2,\dots$ be independent
copies of $X_\sgn$ and $S^\sgn_m=X^\sgn_1+\dots+X^\sgn_m$, with $S^\sgn_0=0$. Then
$S^\sgn_m$ has the law $u^\sgn_m$, and since $X^\sgn_m$ is independent of
$S^\sgn_{m-1}$, for $m,k\ge1$
\begin{equation}\label{eq:q-cover}
 \Prob\bigl(S^\sgn_{m-1}<k\le S^\sgn_m\bigr)=\sum_{r<k}u^\sgn_{m-1}(r)\,T_\sgn(k-r)=q^\sgn_m(k).
\end{equation}
So $q^\sgn_m(k)$ is the probability that the $m$-th run covers level $k$ in the time
of its sign. Every run has length at least $1$, so $S^\sgn_m$ is strictly increasing
in $m$, and for each $k\ge1$ exactly one $m\ge1$ has $S^\sgn_{m-1}<k\le S^\sgn_m$.
Hence, for $k\ge1$,
\begin{equation}\label{eq:kernel-sums}
 \sum_{m\ge1}u^\sgn_m(k)=u_\sgn(k),\qquad
 \sum_{m\ge1}q^\sgn_m(k)=1,\qquad
 \sum_{m\ge1}q^\sgn_{m+1}(k)=1-T_\sgn(k),
\end{equation}
where the third sum is the second minus its term $q^\sgn_1=T_\sgn$. Summing the
middle expression of~\eqref{eq:q-cover} over $m$ gives the covering identity
\begin{equation}\label{eq:cover}
 \sum_{m=1}^{k}T_\sgn(m)\,u_\sgn(k-m)=1\qquad(k\ge1).
\end{equation}
Under (RV) and (SPAN) for the sign $\sgn$ the renewal theorem holds,
\begin{equation}\label{eq:renewal-thm}
 u_\sgn(j)\to1/\mu_\sgn\qquad(j\to\infty),
\end{equation}
and we prove it in local form in Lemma~\ref{lem:index} (see
also~\cite[Section~XIII.3]{feller1968}).

The boundary kernel at the plus end records where a plus run can begin. For a
start of sign $\sgn_0$ let
\[
 \kappa^{(\sgn_0)}_\plus(j)=\Prob\bigl(\text{a plus run starts immediately after the
 $j$-th minus slot}\bigr),
\]
where for $j=0$ the event is that a plus run starts at slot $1$. Count time in
minus slots only. In both cases the ends of the minus runs form a renewal
process with interarrival law $X_\minus$, started at $0$. A plus run starts right
after the $j$-th minus slot exactly when $j$ is a renewal epoch, with one
exception. For a minus start the epoch $j=0$ is followed by the first minus run.
Hence
\begin{equation}\label{eq:kappa}
 \kappa^{(\plus)}_\plus(j)=u_\minus(j)\quad(j\ge0),\qquad
 \kappa^{(\minus)}_\plus(j)=u_\minus(j)\ones_{\{j\ge1\}} .
\end{equation}
The kernel at the minus end is defined in the same way with the signs exchanged.
Every statement below is made at the plus end and holds at the minus end with
the signs exchanged.

We refine $\kappa^{(\sgn_0)}_\plus(i)$ by the number of plus slots before the
run. For $i\ge0$ and $x\ge0$ let $\rho^{(\sgn_0)}_i(x)$ be the probability that a
plus run starts immediately after the $i$-th minus slot (at slot $1$ if $i=0$) and
that exactly $x$ plus slots precede it. At most one $x$ is possible for each
realization, so
\begin{equation}\label{eq:rho-mass}
 \sum_{x\ge0}\rho^{(\sgn_0)}_i(x)=\kappa^{(\sgn_0)}_\plus(i).
\end{equation}
A plus run that starts after the $i$-th minus slot is preceded by $m$ complete
minus runs with $S^\minus_m=i$, and by $m$ plus runs for a plus start and $m-1$
plus runs for a minus start. The two families of runs are independent, so
\begin{equation}\label{eq:rho-rep}
 \rho^{(\plus)}_i(x)=\sum_{m\ge0}u^\minus_m(i)\,u^\plus_m(x),\qquad
 \rho^{(\minus)}_i(x)=\sum_{m\ge1}u^\minus_m(i)\,u^\plus_{m-1}(x).
\end{equation}

\subsection{Exact identities}\label{sec:identities}

The following identities hold for every $N$ and need no hypothesis.

\begin{proposition}[the fresh-plus identity]\label{prop:fresh-identity}
For $k,j\ge1$ with $k+j=N$,
\[
 F^\plus_N(k)=\sum_{m\ge1}\bigl[u^\minus_m(j)\,q^\plus_{m+1}(k)+u^\plus_m(k)\,q^\minus_m(j)\bigr],
\]
and $F^\plus_N(N)=T_\plus(N)$, $F^\plus_N(0)=0$. Exchanging the signs gives the
fresh-minus identity
$F^\minus_N(k)=\sum_m[u^\plus_m(k)q^\minus_{m+1}(j)+u^\minus_m(j)q^\plus_m(k)]$ for
$k,j\ge1$, with $F^\minus_N(0)=T_\minus(N)$ and $F^\minus_N(N)=0$.
\end{proposition}

\begin{proof}
We split according to the run that contains slot $N$. For a fresh plus start the
run lengths $X^\plus_1,X^\plus_2,\dots$ and $X^\minus_1,X^\minus_2,\dots$ are
independent, and the runs alternate starting with a plus run. So the $i$-th plus run
fills the slots in $(S^\plus_{i-1}+S^\minus_{i-1},\,S^\plus_i+S^\minus_{i-1}]$, the
$i$-th minus run fills those in $(S^\plus_i+S^\minus_{i-1},\,S^\plus_i+S^\minus_i]$,
and these intervals partition $\{1,2,\dots\}$. Let $k+j=N$ with $0\le k\le N$.

If slot $N$ lies in the $i$-th plus run, the first $N$ slots contain exactly the
$S^\minus_{i-1}$ minus slots of the first $i-1$ minus runs, so $A_N=k$ means
$S^\minus_{i-1}=j$, and then slot $N=k+j$ lies in that run exactly when
$S^\plus_{i-1}<k\le S^\plus_i$. The two events concern independent sequences, so
by~\eqref{eq:q-cover} this case has probability $u^\minus_{i-1}(j)\,q^\plus_i(k)$. If
slot $N$ lies in the $i$-th minus run, then $A_N=S^\plus_i$, so $A_N=k$ means
$S^\plus_i=k$, and slot $N$ lies in that run exactly when $S^\minus_{i-1}<j\le
S^\minus_i$. This case has probability $u^\plus_i(k)\,q^\minus_i(j)$.

The cases are disjoint. For $k,j\ge1$ the plus case with $i=1$ vanishes because
$u^\minus_0(j)=0$, and $m=i-1$ in the plus case and $m=i$ in the minus case give the
formula. For $k=N$ only the plus case with $i=1$ survives, with probability
$q^\plus_1(N)=T_\plus(N)$. For $k=0$ the plus case is impossible because slot $N$ is
a plus slot, and the minus case is impossible because $S^\plus_i\ge1$. The fresh
minus start is the same with the signs exchanged.
\end{proof}

\begin{proposition}[the first-run decomposition]\label{prop:first-run}
For every preparation and every $0\le k\le N$,
\[
 P_N(k)=Q_\plus(N)\ones_{\{k=N\}}+Q_\minus(N)\ones_{\{k=0\}}
 +\sum_{d<N}\nu_\plus(d)\,F^\minus_{N-d}(k-d)+\sum_{d<N}\nu_\minus(d)\,F^\plus_{N-d}(k).
\]
In particular $P_N(N)=Q_\plus(N)$ and $P_N(0)=Q_\minus(N)$.
\end{proposition}

\begin{proof}
We condition on the first run. If $\Sigma=\plus$ and $D\ge N$, then $A_N=N$, and
this has probability $Q_\plus(N)$. If $\Sigma=\plus$ and $D=d<N$, the first $d$
slots are plus slots and slots $d+1,\dots,N$ are the first $N-d$ slots of a fresh
walk of sign $\minus$, independent of $(\Sigma,D)$. So $A_N=d+A'$, where $A'$ has
the law $F^\minus_{N-d}$. The case $\Sigma=\minus$ is the same with $A_N=0$ when
$D\ge N$ and $A_N$ of law $F^\plus_{N-d}$ when $D=d<N$. This gives the formula. For
$d<N$, Proposition~\ref{prop:fresh-identity} gives $F^\minus_{N-d}(N-d)=0$ and
$F^\plus_{N-d}(0)=0$, and $F^\plus_{N-d}(N)=F^\minus_{N-d}(-d)=0$ trivially, so only
the atoms survive at $k=N$ and $k=0$.
\end{proof}

\begin{lemma}[the Green identity]\label{lem:green}
For each sign and every $b\ge1$,
\[
 \sum_{m\ge1}\Pi^\sgn_m(b)=1+\mu_{\osgn}\,u_\sgn(b).
\]
\end{lemma}

\begin{proof}
Take a fresh walk of sign $\sgn$ continued forever, with runs of sign $\sgn$ of lengths
$X^\sgn_1,X^\sgn_2,\dots$ and runs of sign $\osgn$ of lengths $X^{\osgn}_1,X^{\osgn}_2,\dots$,
all independent, the $i$-th run of sign $\osgn$ coming right after the $i$-th run of
sign $\sgn$. Let $C_m$ be the number of slots of sign $\sgn$ among the first $m$. By
Tonelli's theorem $\sum_m\Pi^\sgn_m(b)=\E\,\#\{m\ge1:C_m=b\}$. Let $t_b<t_{b+1}$ be the
$b$-th and $(b+1)$-th slots of sign $\sgn$. Then $C_m=b$ exactly for $t_b\le m<t_{b+1}$.
If slot $t_b$ is not the last slot of its run, then $t_{b+1}=t_b+1$. If it is, it ends
the $i$-th run of sign $\sgn$ for the unique $i$ with $S^\sgn_i=b$, and then
$t_{b+1}-t_b=1+X^{\osgn}_i$. Hence
$\#\{m:C_m=b\}=1+\sum_{i\ge1}\ones_{\{S^\sgn_i=b\}}X^{\osgn}_i$. Since $X^{\osgn}_i$ is
independent of $S^\sgn_i$, taking expectations and using~\eqref{eq:kernel-sums} gives
$1+\mu_{\osgn}u_\sgn(b)$.
\end{proof}

The next identity isolates the final run, which is the run that contains slot $N$.

\begin{proposition}[the final-run identity]\label{prop:final-run}
For a fresh plus start, $N\ge1$ and $0\le j<N$,
\begin{equation}\label{eq:final-explicit}
\begin{split}
 P_N(N-j)={}&\sum_{r\ge0}u^\minus_r(j)\,\E\bigl[T_\plus(N-j-S^\plus_r);\,S^\plus_r\le N-j-1\bigr]\\
 &+\sum_{m=1}^{j}T_\minus(m)\sum_{r\ge0}u^\minus_r(j-m)\,u^\plus_{r+1}(N-j).
\end{split}
\end{equation}
The first sum is the event that the final run is a plus run, and the second the
event that it is a minus run of length at least $m$ that starts at slot $N-m+1$. Only
$r\le j$ contributes, and for $j=0$ the identity reads $P_N(N)=T_\plus(N)$.
\end{proposition}

\begin{proof}
We use the notation of the proof of Proposition~\ref{prop:fresh-identity}. The
$(r+1)$-th plus run starts at slot $S^\plus_r+S^\minus_r+1$ with $S^\plus_r$ plus slots
before it. It is the final run and $A_N=N-j$ exactly when the slots before it contain
$j$ minus slots, that is $S^\minus_r=j$, and its length satisfies
$X^\plus_{r+1}\ge N-j-S^\plus_r\ge1$. These conditions involve independent variables,
which gives the term $r$ of the first sum. The $(r+1)$-th minus run starts at slot
$S^\plus_{r+1}+S^\minus_r+1$ with $S^\plus_{r+1}$ plus slots before it. If it is the
final run and starts at slot $N-m+1$, then $A_N=N-j$ means $S^\plus_{r+1}=N-j$, hence
$S^\minus_r=N-m-(N-j)=j-m$, which needs $1\le m\le j$, and its length satisfies
$X^\minus_{r+1}\ge m$. This gives the second sum. Exactly one run is final, so the
cases are disjoint. Finally $u^\minus_r(i)=0$ for $r>i$ and $u^\minus_r(0)=\ones_{\{r=0\}}$.
\end{proof}

By~\eqref{eq:rho-rep} the first sum of~\eqref{eq:final-explicit} equals
$\sum_{x\le N-j-1}\rho^{(\plus)}_j(x)T_\plus(N-j-x)$. Section~\ref{sec:verify} describes
how these identities compute the exact law.

\begin{lemma}[the counting lemma]\label{lem:counting}
Fix $j\ge0$ and $N>j$. On the event $\{A_N=N-j\}$ at most $j+1$ plus runs meet the
first $N$ slots (at most $j$ if the first run is a minus run), and at least one of
them has at least $(N-j)/(j+1)$ slots among the first $N$. For a fresh start of
either sign and every $L\ge1$, the probability that $A_N=N-j$ and two plus runs
that meet the first $N$ slots have length at least $L$ is at most
$(j+1)^2T_\plus(L)^2$.
\end{lemma}

\begin{proof}
On $\{A_N=N-j\}$ the first $N$ slots contain $j$ minus slots. Let $r$ be the number of
plus runs that meet them. Between two consecutive such runs lies a minus run inside
the first $N$ slots, so $r-1\le j$, and if the first run is a minus run it adds one
more minus slot, so $r\le j$. The $N-j$ plus slots are shared among $r\le j+1$ runs,
so one of them has at least $(N-j)/(j+1)$ slots. For a fresh start the plus runs have
independent lengths $X^\plus_1,X^\plus_2,\dots$, and on the event of the lemma there
are $1\le a<b\le j+1$ with $X^\plus_a\ge L$ and $X^\plus_b\ge L$. The union bound gives
at most $\binom{j+1}2T_\plus(L)^2\le(j+1)^2T_\plus(L)^2$, and counting only the slots
among the first $N$ makes the event smaller.
\end{proof}
